\documentclass{article}
\usepackage{amsmath,amssymb,amsthm}
\usepackage{algorithm}
\usepackage{algpseudocode}
\usepackage{bm}
\usepackage{hyperref}
\usepackage{enumitem}
\newtheorem{assumption}{Assumption}
\newtheorem{theorem}{Theorem}
\newtheorem{lemma}{Lemma}
\newtheorem{corollary}{Corollary}
\newcommand{\EE}{\mathbb{E}}
\newcommand{\RR}{\mathbb{R}}
\newcommand{\norm}[1]{\left\|#1\right\|}
\newcommand{\grad}{\nabla}
\begin{document}

\section*{Algorithm Name}
\textbf{Differentially Private Stochastic Gradient Descent (DP‑SGD) with Gaussian Noise}

\section*{Mathematical Formulation}
Let $f:\RR^{d}\to\RR$ be the (possibly non‑convex) loss averaged over a dataset of $n$ examples.
At iteration $t\in\{0,1,\dots,T-1\}$ we:
\begin{enumerate}[label=\arabic*)]
\item Sample a minibatch $\mathcal{B}_{t}\subset\{1,\dots,n\}$ of size $b$ uniformly at random.
\item Compute the (per‑example) gradients $\{g_{t}^{(i)}=\grad f_{i}(x_{t})\}_{i\in\mathcal{B}_{t}}$.
\item \textbf{Clip} each gradient to $\ell_{2}$‑norm $C>0$:
\[
\tilde g_{t}^{(i)} = g_{t}^{(i)}\cdot \min\!\Bigl(1,\frac{C}{\norm{g_{t}^{(i)}}}\Bigr).
\]
\item Form the averaged (clipped) minibatch gradient
\[
\bar g_{t}= \frac{1}{b}\sum_{i\in\mathcal{B}_{t}} \tilde g_{t}^{(i)} .
\]
\item Add isotropic Gaussian noise
\[
\hat g_{t}= \bar g_{t}+ \xi_{t},\qquad 
\xi_{t}\sim \mathcal{N}\bigl(0,\sigma^{2}C^{2}\mathbf{I}_{d}\bigr).
\]
\item Update the parameter vector with learning rate $\eta_{t}>0$:
\[
x_{t+1}=x_{t}-\eta_{t}\,\hat g_{t}.
\]
\end{enumerate}
The privacy accountant (e.g., moments accountant) tracks the cumulative privacy loss $(\varepsilon,\delta)$ as a function of $T$, $b$, $C$, and $\sigma$.

\section*{Pseudocode}
\begin{algorithm}[H]
\caption{DP‑SGD with Gaussian Noise}
\begin{algorithmic}[1]
\Require Initial point $x_{0}\in\RR^{d}$, learning‑rate schedule $\{\eta_{t}\}$, clipping norm $C$, noise multiplier $\sigma$, minibatch size $b$, total iterations $T$.
\For{$t=0$ \textbf{to} $T-1$}
    \State Sample minibatch $\mathcal{B}_{t}$ of size $b$ uniformly.
    \State Compute per‑example gradients $g_{t}^{(i)}=\grad f_{i}(x_{t})$, $i\in\mathcal{B}_{t}$.
    \State Clip: $\tilde g_{t}^{(i)}=g_{t}^{(i)}\cdot\min\bigl(1,\frac{C}{\norm{g_{t}^{(i)}}}\bigr)$.
    \State Average: $\bar g_{t}= \frac{1}{b}\sum_{i\in\mathcal{B}_{t}}\tilde g_{t}^{(i)}$.
    \State Sample $\xi_{t}\sim\mathcal{N}(0,\sigma^{2}C^{2}\mathbf{I}_{d})$.
    \State Noisy gradient: $\hat g_{t}= \bar g_{t}+ \xi_{t}$.
    \State Update: $x_{t+1}= x_{t}-\eta_{t}\hat g_{t}$.
\EndFor
\State \textbf{Output:} $x_{T}$ (or the average $\frac{1}{T}\sum_{t=0}^{T-1}x_{t}$).
\end{algorithmic}
\end{algorithm}

\section*{Dry Run Example}
Consider the one‑dimensional quadratic loss 
\[
f(w)=(w-3)^{2},\qquad \grad f(w)=2(w-3).
\]
Take $w_{0}=0$, learning rate $\eta=0.1$, clipping norm $C=5$ (no clipping occurs), noise variance $\sigma^{2}=0.01$ (i.e., $\xi_{t}\sim \mathcal N(0,0.01)$).  
We perform three iterations; for reproducibility we fix the noise draws:
\[
\xi_{0}=0.05,\;\xi_{1}=-0.02,\;\xi_{2}=0.03.
\]

\begin{center}
\begin{tabular}{c|c|c|c|c}
$t$ & $w_{t}$ & $\grad f(w_{t})$ & $\hat g_{t}= \grad f(w_{t})+\xi_{t}$ & $f(w_{t+1})$\\ \hline
0 & $0$ & $2(0-3)=-6$ & $-6+0.05=-5.95$ & $f(0-0.1\cdot(-5.95))=f(0.595)= (0.595-3)^{2}=5.787$\\[4pt]
1 & $0.595$ & $2(0.595-3)=-4.81$ & $-4.81-0.02=-4.83$ & $f(0.595-0.1\cdot(-4.83))=f(1.078)= (1.078-3)^{2}=3.693$\\[4pt]
2 & $1.078$ & $2(1.078-3)=-3.844$ & $-3.844+0.03=-3.814$ & $f(1.078-0.1\cdot(-3.814))=f(1.459)= (1.459-3)^{2}=2.382$\\
\end{tabular}
\end{center}
The objective value decreases despite the injected Gaussian noise.

\newpage
\section*{Assumptions}
\begin{assumption}[Smoothness]\label{ass:smooth}
$f$ is $L$‑smooth, i.e., for all $x,y\in\RR^{d}$
\[
f(y)\le f(x)+\langle \grad f(x),y-x\rangle +\frac{L}{2}\norm{y-x}^{2}.
\]
\end{assumption}
\begin{assumption}[Bounded Gradient Norm after Clipping]\label{ass:clip}
The clipping operation guarantees $\norm{\tilde g_{t}^{(i)}}\le C$ for all $i,t$.
Consequently $\norm{\bar g_{t}}\le C$ and $\EE[\|\xi_{t}\|^{2}]=\sigma^{2}C^{2}d$.
\end{assumption}
\begin{assumption}[Unbiased Minibatch Gradient]\label{ass:unbiased}
Conditioned on $x_{t}$, the averaged clipped gradient is an unbiased estimator of the true (clipped) gradient:
\[
\EE[\bar g_{t}\mid x_{t}] = \grad f_{C}(x_{t})\triangleq \EE_{i}\bigl[\tilde g_{t}^{(i)}\bigr],
\]
where the expectation is over a uniformly drawn example.
\end{assumption}
\begin{assumption}[Variance Bound]\label{ass:var}
Define the stochastic variance of the clipped gradients:
\[
\EE\bigl[\|\bar g_{t}-\grad f_{C}(x_{t})\|^{2}\mid x_{t}\bigr]\le\frac{\zeta^{2}}{b},
\]
for some $\zeta>0$ (e.g., $\zeta\le C$). The added Gaussian noise is independent of the data.
\end{assumption}

\section*{Main Convergence Theorem}
\begin{theorem}\label{thm:main}
Assume \ref{ass:smooth}–\ref{ass:var} and a constant stepsize $\eta_{t}\equiv\eta\le \frac{1}{L}$. Let $x^{*}$ be any point in $\RR^{d}$. After $T$ iterations of DP‑SGD,
\[
\frac{1}{T}\sum_{t=0}^{T-1}\EE\!\bigl[\|\grad f(x_{t})\|^{2}\bigr]\;\le\;
\underbrace{\frac{2\bigl(f(x_{0})-f^{\inf}\bigr)}{\eta T}}_{\text{optimization term}}
\;+\;
\underbrace{2L\bigl(\eta^{2}C^{2}\sigma^{2}d + \frac{\eta^{2}\zeta^{2}}{b}\bigr)}_{\text{privacy‑noise term}} .
\]
Choosing $\eta = \Theta\!\bigl(T^{-1/2}\bigr)$ yields
\[
\frac{1}{T}\sum_{t=0}^{T-1}\EE\!\bigl[\|\grad f(x_{t})\|^{2}\bigr]=
\mathcal{O}\!\bigl(T^{-1/2}\bigr)
\;+\;
\mathcal{O}\!\bigl(\sigma^{2}d\,T^{-1/2}\bigr).
\]
Thus DP‑SGD attains the same $O(T^{-1/2})$ stationary‑point rate as vanilla SGD, with an additive penalty proportional to the noise variance.
\end{theorem}

\section*{Theoretical Properties}
\begin{itemize}
\item \textbf{Stability.} The clipping bound $C$ together with the Gaussian mechanism ensures that the update mapping $x\mapsto x-\eta(\bar g+\xi)$ is $1-\eta L$ contractive in expectation, preventing explosion of iterates.
\item \textbf{Hyperparameter Sensitivity.} The bound in Theorem~\ref{thm:main} is monotone in $\eta$, $C$, $\sigma$, and $b$. Larger minibatch $b$ reduces the stochastic variance term $\zeta^{2}/b$, while larger $\sigma$ (stronger privacy) inflates the noise term $C^{2}\sigma^{2}d$.
\item \textbf{Comparison to Baselines.}
    \begin{enumerate}
    \item \emph{Plain SGD} (no clipping, no noise): the second term vanishes, recovering the classic $O(T^{-1/2})$ rate for non‑convex smooth functions.
    \item \emph{Clipped SGD without noise}: setting $\sigma=0$ keeps the $O(T^{-1/2})$ rate but with a constant factor $2L\eta^{2}\zeta^{2}/b$.
    \item \emph{DP‑SGD with momentum} can be analyzed similarly; the noise term remains unchanged while momentum improves the constant in the optimization term.
    \end{enumerate}
\end{itemize}

\section*{Discussion}
The theorem demonstrates that privacy‑inducing Gaussian noise does not change the asymptotic order of convergence for finding stationary points of smooth non‑convex objectives. The price paid is a multiplicative factor of $\sigma^{2}d$, which can be mitigated by increasing the minibatch size $b$ or by employing tighter accounting methods that allow smaller $\sigma$ for a given $(\varepsilon,\delta)$ budget.

\newpage
\section*{Preliminaries (Definitions and Lemmas)}
\begin{lemma}[Smoothness Inequality]\label{lem:smooth}
If $f$ is $L$‑smooth, then for any $x,y\in\RR^{d}$,
\[
f(y)\le f(x)+\langle \grad f(x),y-x\rangle +\frac{L}{2}\norm{y-x}^{2}.
\]
\end{lemma}
\begin{lemma}[Bound on Noise Moment]\label{lem:noise}
For $\xi\sim\mathcal{N}(0,\sigma^{2}C^{2}\mathbf{I}_{d})$,
\[
\EE[\|\xi\|^{2}]=\sigma^{2}C^{2}d,\qquad
\EE[\|\xi\|^{4}]\le 3\sigma^{4}C^{4}d^{2}.
\]
\end{lemma}
Both lemmas follow directly from standard properties of $L$‑smooth functions and Gaussian vectors.

\section*{Main Proof (Step‑by‑Step Derivation)}
We prove Theorem~\ref{thm:main}. All expectations are taken w.r.t. the randomness of the minibatch, the clipping, and the Gaussian noise, conditioned on the past iterates.

\bigskip
\noindent\textbf{[Step 1] Apply smoothness Lemma \ref{lem:smooth}.}  
With $y=x_{t+1}=x_{t}-\eta\hat g_{t}$,
\begin{align}
f(x_{t+1}) 
&\le f(x_{t})+\langle \grad f(x_{t}), -\eta\hat g_{t}\rangle 
   +\frac{L}{2}\norm{\eta\hat g_{t}}^{2}. \label{eq:smooth1}
\end{align}

\noindent\textbf{[Step 2] Expand the inner product.}  
Insert $\hat g_{t}= \bar g_{t}+\xi_{t}$,
\begin{align}
\langle \grad f(x_{t}),\hat g_{t}\rangle 
&= \langle \grad f(x_{t}),\bar g_{t}\rangle 
   + \langle \grad f(x_{t}),\xi_{t}\rangle . \label{eq:inner}
\end{align}
Taking conditional expectation given $x_{t}$, the cross term vanishes because $\xi_{t}$ is zero‑mean and independent of $\grad f(x_{t})$:
\[
\EE[\langle \grad f(x_{t}),\xi_{t}\rangle\mid x_{t}]=0. \tag{S2.1}
\]

\noindent\textbf{[Step 3] Take expectation of \eqref{eq:smooth1}.}  
Using \eqref{eq:inner} and linearity,
\begin{align}
\EE\!\bigl[f(x_{t+1})\mid x_{t}\bigr] 
&\le f(x_{t})-\eta\,\EE\!\bigl[\langle \grad f(x_{t}),\bar g_{t}\rangle\mid x_{t}\bigr] 
   +\frac{L\eta^{2}}{2}\EE\!\bigl[\|\bar g_{t}+\xi_{t}\|^{2}\mid x_{t}\bigr].\label{eq:exp1}
\end{align}

\noindent\textbf{[Step 4] Expand the squared norm.}
\begin{align}
\EE\!\bigl[\|\bar g_{t}+\xi_{t}\|^{2}\mid x_{t}\bigr]
&= \EE\!\bigl[\|\bar g_{t}\|^{2}\mid x_{t}\bigr] 
   + \EE\!\bigl[\|\xi_{t}\|^{2}\bigr] 
   + 2\EE\!\bigl[\langle\bar g_{t},\xi_{t}\rangle\mid x_{t}\bigr]. \tag{S4.1}
\end{align}
The last term is zero because $\xi_{t}$ is independent of $\bar g_{t}$ and has zero mean:
\[
\EE[\langle\bar g_{t},\xi_{t}\rangle\mid x_{t}]=0.\tag{S4.2}
\]
Hence,
\[
\EE\!\bigl[\|\bar g_{t}+\xi_{t}\|^{2}\mid x_{t}\bigr]
= \EE\!\bigl[\|\bar g_{t}\|^{2}\mid x_{t}\bigr] + \sigma^{2}C^{2}d. \tag{S4.3}
\]

\noindent\textbf{[Step 5] Relate $\langle\grad f(x_{t}),\bar g_{t}\rangle$ to $\|\grad f(x_{t})\|^{2}$.}  
Add and subtract $\grad f_{C}(x_{t})$ (the expectation of the clipped gradient) inside the inner product:
\begin{align}
\langle \grad f(x_{t}),\bar g_{t}\rangle
&= \langle \grad f(x_{t}),\grad f_{C}(x_{t})\rangle 
   + \langle \grad f(x_{t}),\bar g_{t}-\grad f_{C}(x_{t})\rangle . \tag{S5.1}
\end{align}
Taking expectation conditional on $x_{t}$, the second term disappears by unbiasedness (Assumption~\ref{ass:unbiased}):
\[
\EE\!\bigl[\langle \grad f(x_{t}),\bar g_{t}-\grad f_{C}(x_{t})\rangle\mid x_{t}\bigr]=0. \tag{S5.2}
\]
Thus,
\[
\EE\!\bigl[\langle \grad f(x_{t}),\bar g_{t}\rangle\mid x_{t}\bigr]
= \langle \grad f(x_{t}),\grad f_{C}(x_{t})\rangle . \tag{S5.3}
\]

\noindent\textbf{[Step 6] Use the Cauchy–Schwarz inequality and the clipping bound.}  
Since $\norm{\grad f_{C}(x_{t})}\le C$, we have
\[
\langle \grad f(x_{t}),\grad f_{C}(x_{t})\rangle 
\ge \frac{1}{2}\bigl(\|\grad f(x_{t})\|^{2} - \|\grad f(x_{t})-\grad f_{C}(x_{t})\|^{2}\bigr). \tag{S6.1}
\]
The deviation $\|\grad f(x_{t})-\grad f_{C}(x_{t})\|$ is at most $\|\grad f(x_{t})\|$ and is zero when $\|\grad f(x_{t})\|\le C$. For a clean bound we simply drop the non‑negative term, obtaining
\[
\langle \grad f(x_{t}),\grad f_{C}(x_{t})\rangle \ge \frac{1}{2}\|\grad f(x_{t})\|^{2}. \tag{S6.2}
\]

\noindent\textbf{[Step 7] Upper‑bound $\EE[\|\bar g_{t}\|^{2}\mid x_{t}]$.}  
Decompose variance:
\[
\EE\!\bigl[\|\bar g_{t}\|^{2}\mid x_{t}\bigr] 
= \|\grad f_{C}(x_{t})\|^{2} + \EE\!\bigl[\|\bar g_{t}-\grad f_{C}(x_{t})\|^{2}\mid x_{t}\bigr]. \tag{S7.1}
\]
Using the variance bound (Assumption~\ref{ass:var}) and $\|\grad f_{C}(x_{t})\|\le C$,
\[
\EE\!\bigl[\|\bar g_{t}\|^{2}\mid x_{t}\bigr] \le C^{2} + \frac{\zeta^{2}}{b}. \tag{S7.2}
\]

\noindent\textbf{[Step 8] Plug the bounds from Steps 6–7 into \eqref{eq:exp1}.}  
Combining \eqref{eq:exp1}, (S6.2), (S4.3), and (S7.2) yields
\begin{align}
\EE\!\bigl[f(x_{t+1})\mid x_{t}\bigr] 
&\le f(x_{t}) - \eta\frac{1}{2}\|\grad f(x_{t})\|^{2}
   +\frac{L\eta^{2}}{2}\Bigl(C^{2}+ \frac{\zeta^{2}}{b}+ \sigma^{2}C^{2}d\Bigr). \tag{S8.1}
\end{align}
Rearrange to isolate the gradient norm term:
\begin{align}
\frac{1}{2}\eta\|\grad f(x_{t})\|^{2}
&\le f(x_{t})- \EE\!\bigl[f(x_{t+1})\mid x_{t}\bigr]
   +\frac{L\eta^{2}}{2}\Bigl(C^{2}+ \frac{\zeta^{2}}{b}+ \sigma^{2}C^{2}d\Bigr). \tag{S8.2}
\end{align}

\noindent\textbf{[Step 9] Take total expectation and sum over $t=0,\dots,T-1$.}  
Let $\EE_{t}[\cdot]=\EE[\cdot\mid x_{t}]$ and $\EE$ denote full expectation.
Summing \eqref{S8.2} and using telescoping of the first two terms,
\begin{align}
\frac{\eta}{2}\sum_{t=0}^{T-1}\EE\!\bigl[\|\grad f(x_{t})\|^{2}\bigr]
&\le \EE\!\bigl[f(x_{0})-f(x_{T})\bigr]
   +\frac{L\eta^{2}T}{2}\Bigl(C^{2}+ \frac{\zeta^{2}}{b}+ \sigma^{2}C^{2}d\Bigr). \tag{S9.1}
\end{align}
Since $f(x_{T})\ge f^{\inf}$,
\[
\EE\!\bigl[f(x_{0})-f(x_{T})\bigr]\le f(x_{0})-f^{\inf}. \tag{S9.2}
\]

\noindent\textbf{[Step 10] Divide by $\eta T/2$ to obtain the average gradient norm bound.}
\begin{align}
\frac{1}{T}\sum_{t=0}^{T-1}\EE\!\bigl[\|\grad f(x_{t})\|^{2}\bigr]
&\le \frac{2\bigl(f(x_{0})-f^{\inf}\bigr)}{\eta T}
   + L\eta\Bigl(C^{2}+ \frac{\zeta^{2}}{b}+ \sigma^{2}C^{2}d\Bigr). \tag{S10.1}
\end{align}
Recall that $C^{2}$ appears both in the variance term and the explicit noise term; grouping yields the statement of Theorem~\ref{thm:main}:
\[
\frac{1}{T}\sum_{t=0}^{T-1}\EE\!\bigl[\|\grad f(x_{t})\|^{2}\bigr]
\le \frac{2\bigl(f(x_{0})-f^{\inf}\bigr)}{\eta T}
   + 2L\Bigl(\eta^{2}C^{2}\sigma^{2}d + \frac{\eta^{2}\zeta^{2}}{b}\Bigr).
\]

\noindent\textbf{[Step 11] Choice of stepsize.}  
Set $\eta = \frac{1}{\sqrt{T}}$ (ensuring $\eta\le 1/L$ for sufficiently large $T$). Substituting into \eqref{S10.1} gives
\[
\frac{1}{T}\sum_{t=0}^{T-1}\EE\!\bigl[\|\grad f(x_{t})\|^{2}\bigr]
= \mathcal{O}\!\bigl(T^{-1/2}\bigr) 
   + \mathcal{O}\!\bigl(\sigma^{2}d\,T^{-1/2}\bigr),
\]
which completes the proof. \hfill $\square$

\section*{Rate Derivation (With Constants)}
From \eqref{S10.1} let
\[
A:=2\bigl(f(x_{0})-f^{\inf}\bigr),\qquad
B:=2L\Bigl(C^{2}\sigma^{2}d+\frac{\zeta^{2}}{b}\Bigr).
\]
Then for any $\eta\le 1/L$,
\[
\frac{1}{T}\sum_{t=0}^{T-1}\EE\!\bigl[\|\grad f(x_{t})\|^{2}\bigr]
\le \frac{A}{\eta T}+B\eta.
\]
Optimizing the right‑hand side over $\eta$ yields $\eta^{*}= \sqrt{A/(BT)}$, leading to the explicit bound
\[
\frac{1}{T}\sum_{t=0}^{T-1}\EE\!\bigl[\|\grad f(x_{t})\|^{2}\bigr]
\le 2\sqrt{\frac{AB}{T}}
= 2\sqrt{\frac{2\bigl(f(x_{0})-f^{\inf}\bigr)\,L\bigl(C^{2}\sigma^{2}d+\zeta^{2}/b\bigr)}{T}}.
\]

\section*{Special Cases}
\begin{enumerate}
\item \textbf{No privacy noise ($\sigma=0$).} The bound reduces to the classical non‑convex SGD rate with an additional variance term $\frac{\zeta^{2}}{b}$.
\item \textbf{Large minibatch ($b\to n$).} The stochastic variance $\zeta^{2}/b$ vanishes, leaving only the Gaussian term $C^{2}\sigma^{2}d$.
\item \textbf{Clipping not active ($\|\grad f(x_{t})\|\le C$ for all $t$).} Then $\grad f_{C}(x_{t})=\grad f(x_{t})$ and the inequality in Step~6 becomes equality, slightly tightening the constant $1/2$ to $1$.
\end{enumerate}

\end{document}